\documentclass[12pt]{article}

% === Packages ===
\usepackage[margin=1in]{geometry}
\usepackage{amsmath, amssymb}
\usepackage{setspace}
\usepackage{titlesec}
\usepackage{enumitem}
\usepackage{helvet}
\renewcommand{\familydefault}{\sfdefault} % modern sans-serif look

% === Formatting tweaks ===
\setstretch{1.15}
\setlist[enumerate]{itemsep=0.4em, topsep=0.2em}
\titleformat{\section}{\Large\bfseries\sffamily}{\thesection.}{1em}{}
\titleformat{\subsection}{\large\bfseries\sffamily}{\thesubsection}{1em}{}

% === Title ===
\title{\vspace{-2em}\textbf{Applied Mathematics Oral Exam Questions and Model Answers}\\
\large From Euler to Navier--Stokes}
\author{Prepared for Graduate Oral Examination}
\date{}

\begin{document}
\section*{Sample Oral Exam Questions Based on Non-Dimensionalization}
\subsection*{Questions and Model Answers}

% -------------------------------------------
\subsection*{1. Why do we nondimensionalize the Navier--Stokes equations?}

\textbf{Answer.}
Nondimensionalization allows us to:
\begin{itemize}
    \item compare flows of different physical scales,
    \item reduce the number of free parameters,
    \item identify the dominant physical balances,
    \item ensure numerical stability by revealing appropriate scaling.
\end{itemize}
According to the tutorial (see p.~2 of the PDF), you either work in real units
or report dimensionless parameters. Dimensionless variables like
$x' = x/L$, $u'=u/U$, and $t'=tU/L$ appear directly in the IBFE formulation
whenever the fluid equations are nondimensionalized.


% -------------------------------------------
\subsection*{2. Derive the nondimensional incompressible Navier--Stokes equations.}

\textbf{Answer.}
Start from dimensional Navier--Stokes:
\[
\rho\left(
\frac{\partial \mathbf{u}}{\partial t} + (\mathbf{u}\cdot\nabla)\mathbf{u}
\right)
= -\nabla p + \mu\nabla^2\mathbf{u}.
\]

Use the definitions (PDF p.~4):
\[
x'=\frac{x}{L},\quad t'=\frac{tU}{L},\quad u'=\frac{u}{U},
\quad p'=\frac{p}{\rho U^2},\quad \nabla'=L\nabla.
\]
Substituting and rearranging yields
\[
\frac{\partial \mathbf{u}'}{\partial t'}
+ \mathbf{u}'\cdot\nabla'\mathbf{u}'
= -\nabla' p' + \frac{1}{Re}\nabla'^2\mathbf{u}',
\]
where
\[
Re = \frac{\rho UL}{\mu}.
\]


% -------------------------------------------
\subsection*{3. What is the physical meaning of the Reynolds number in IBFE simulations?}

\textbf{Answer.}
From the definition in the tutorial (p.~4), $Re=\rho UL/\mu$ measures the ratio of
inertial to viscous forces.  
In IBFE:
\begin{itemize}
    \item Low $Re$ implies the flow is dominated by viscosity, requiring smaller 
          time steps and possibly implicit treatment of stiffness.
    \item High $Re$ implies inertia-dominated flow, requiring careful resolution
          of vortices and wake structures.
    \item The choice of $Re$ determines how elastic boundary forces scale and how
          strong spreading forces must be on the Eulerian grid.
\end{itemize}
The plot on p.~5 shows how real biological and engineered systems span several
orders of magnitude in $Re$.
% -------------------------------------------
\subsection*{4. How do you nondimensionalize force per unit length in IBFE?}

\textbf{Answer.}
If a Lagrangian fiber applies a force per unit length $F$, the dimensionless force
depends on the Reynolds regime (p.~7):
\[
F' = 
\begin{cases}
F / (\rho L U^2), & Re>1,\\[6pt]
F / (\mu U), & Re<1.
\end{cases}
\]
This scaling ensures consistency with the nondimensional Navier--Stokes
equations so that structural forcing terms are comparable in magnitude to the
fluid terms being solved.

% -------------------------------------------
\subsection*{5. How is the dimensionless beam bending stiffness $k'_{\text{beam}}$ derived?}

\textbf{Answer.}
The bending stiffness in 2D has units of $\mathrm{N\cdot m}$ (PDF p.~8).
For $Re>1$, the nondimensional bending stiffness is (p.~12):
\[
k'_{\text{beam}}=\frac{k_{\text{beam}}}{\rho U^2 L^3}.
\]
This comes from comparing the physical bending force scale
($k_{\text{beam}}/L^2$) to the inertial fluid stress scale
($\rho U^2$).  
Similarly, for $Re<1$:
\[
k'_{\text{beam}}=\frac{k_{\text{beam}}}{\mu U L^2}.
\]

% -------------------------------------------
\subsection*{6. Estimate bending stiffness and explain how it affects numerical stability.}

\textbf{Answer.}
Page~15 notes that $k'_{\text{beam}}\approx 10$ produces moderate bending.
In IBFE:
\begin{itemize}
    \item Too large $k_{\text{beam}}$ leads to numerical stiffness and requires
          smaller time steps or implicit treatment.
    \item Too small $k_{\text{beam}}$ produces exaggerated deformation and mesh
          distortion.
    \item Estimating $k_{\text{beam}}$ using physical $U_{\max}$ and boundary 
          length ensures meaningful nondimensionalization.
\end{itemize}

% -------------------------------------------
\subsection*{7. How do you estimate a spring constant for immersed boundary structures?}

\textbf{Answer.}
From p.~16,
\[
k_{\text{spring}} \approx \frac{F_b}{x_s},
\]
where $x_s$ is the desired displacement and $F_b$ depends on Reynolds number:
\[
F_b \sim
\begin{cases}
\frac{1}{2}\rho\,dx^{(\mathrm{dim}-1)}U^2, & Re>1,\\[6pt]
\mu\,dx\,U, & Re<1.
\end{cases}
\]
In IBFE, this determines how stiff the immersed material is relative to fluid
stresses, preventing unrealistic penetration or stretching of elastic boundaries.

% -------------------------------------------
\subsection*{8. What dimensionless groups govern oscillatory or pumping flows?}

\textbf{Answer.}
Several groups appear in your notes:
\begin{itemize}
    \item \textbf{Dimensionless time:} $t'=tU/L$ (p.~6).
    \item \textbf{Dimensionless frequency:} $f'=fL/U$ (p.~6).
    \item \textbf{Strouhal number:} $St = fA/U$ (pp.~18--19),
          governing vortex shedding and propulsive efficiency
          ($0.2<St<0.4$ for efficient swimmers).
    \item \textbf{Womersley number:} $\alpha = L\sqrt{\omega/\nu}$ (p.~20),
          characterizing oscillatory flow in biological systems such as arteries.
\end{itemize}
These determine how oscillations couple to fluid inertia and viscosity, crucial
in pumping or swimming IBFE simulations.

% -------------------------------------------
\subsection*{9. How does nondimensionalization guide mesh and time-step choices?}

\textbf{Answer.}
The dimensionless groups determine:
\begin{itemize}
    \item the smallest spatial scales needing resolution 
          (e.g., boundary layers for large $Re$ or Womersley number),
    \item the maximum allowable Courant number for explicit schemes,
    \item stiffness in structural models for large $k'_{\text{beam}}$ or 
          $k'_{\text{spring}}$,
    \item time-step requirements for coupled fluid–structure interaction.
\end{itemize}
For example, high Strouhal number flows induce rapid vortex shedding,
necessitating careful time-step selection.  
Similarly, low-Re IBFE simulations require fine Eulerian grids because viscous 
diffusion dominates.

% -------------------------------------------
\subsection*{10. Explain the physical meaning of the Péclet number and how it applies in IBFE.}

\textbf{Answer.}
From p.~22,
\[
Pe = \frac{\text{advective transport rate}}
         {\text{diffusive transport rate}}
       = \frac{LU}{D}.
\]
In IBFE:
\begin{itemize}
    \item For passive scalars or mass/chemical transport, $Pe$ determines
          whether advection or diffusion dominates.
    \item If $Pe \gg 1$, steep gradients require enhanced mesh resolution or
          stabilization.
    \item If $Pe \ll 1$, diffusion smooths fields, potentially relaxing mesh
          requirements.
\end{itemize}
This impacts numerical performance whenever diffusion terms are included in
the Eulerian equations solved alongside fluid–structure interaction.

\end{document}